% Template for post or presentation abstracts submitted to ILAS 2022
% 1. Fill in the presenter's information (name, affiliation, email
% address) and talk information (title of presentation, abstract).
% 2. Compile to make sure there are no errors.
% 3. Save the .tex file for your abstract with a distinctive name,
% ideally "FamilynameGivenname.tex".
% 4. Upload your .tex file to https://clr.ie/129557
%
% * Please keep your LaTeX commands simple, and avoid the use of
%    user-defined macros.
% * Include details of co-authors and funding acknowledgements after the abstract.
% * Upload only the LaTeX file (with .tex extension).
% * Questions: Contact Niall Madden (Niall.Madden@NUIGalway.ie)

\documentclass[a4paper]{amsart}
\usepackage{amssymb}
\usepackage{hyperref}
\pagestyle{empty}
\date{}

%% IGNORE THE NEXT 4 LINES
\makeatletter %
\def\labelenumi{\theenumi.}
\def\theenumi{\@arabic\c@enumi}
\makeatother
%% STOP IGNORING NOW

\begin{document}

\title{The $H$-join of arbitrary families of graphs - the universal adjacency spectrum}
\author{Helena Gomes} %Presenting author only
\address{Center for Research and Development in Mathematics and Applications - CIDMA, University of Aveiro} % Affiliation of presenting author
\email{hgomes@ua.pt} % Email address of presenting author only

\maketitle

% The abstract  ***no more than one page***

The $H$-join of a family of graphs $\mathcal{G}=\{G_1, \dots, G_p\}$, also called generalized composition, $H[G_1, \dots, G_p]$, where all graphs are undirected, simple and finite, is the graph obtained by replacing each vertex $i$ of $H$ by $G_i$ and adding to the edges of all graphs in $\mathcal{G}$ the edges of the join $G_i \vee G_j$, for every edge $ij$ of $H$. Some well known graph operations are particular cases of the $H$-join of a family of graphs $\mathcal{G}$ as it is the case of the lexicographic product (also called composition) of two graphs $H$ and $G$, $H[G]$. During long time the known expressions for the determination of the entire spectrum of the $H$-join in terms of the spectra of its components (that is, graphs in $\mathcal{G}$) and an associated matrix, related with the main eigenvalues of the components and the graph $H$, were limited to families $\mathcal{G}$ of regular graphs. In this work, with an approach based on the walk-matrix, we extend such a determination, as well as the determination of the characteristic polynomial, to the universal adjacency matrix of the $H$-join of families of arbitrary graphs. From the obtained results, the eigenvectors of the universal adjacency matrix of the $H$-join can also be determined in terms of some eigenvectors of the universal adjacency matrices of the components and an associated matrix.%
%Please, follow the instructions at \url{http://ilas2022.ie/abstracts/}
%for submitting your abstract. Make sure you upload the \texttt{.tex}
%file, and not any others.
%
%Try to keep your file as simple as possible, so that it is easily
%rendered in various formats. Give the file a name that readily
%identifies it with you, for example \texttt{OlgaTaussky.tex}.
%
%If you wish to include citations, you can do so like this:
%\cite{my_key_for_OT49}. Use a bibitem key that is likely be be unique to
%you (e.g., don't use the default Math Reviews key).
%Any uncited bibliography entries will also
%appear. If you have none, then delete that section.
%
% Coauthor details here (optional - delete if not applicable)
% and funding acknowledgment (optional - delete if not applicable)

\bigskip

\emph{This is joint work with Domingos M. Cardoso (University of Aveiro) and Sofia Pinheiro (University of Aveiro).
Supported by the Center for Research and Development in Mathematics and Applications (CIDMA) which is financed by national funds through Funda\c{c}\~{a}o para a Ci\^{e}ncia e a Tecnologia (FCT), Grant UIDB/04106/2020.}
%
%\begin{thebibliography}{1}
%\bibitem{my_key_for_OT49}
%Olga Taussky.
%\newblock A recurring theorem on determinants.
%\newblock {\em  Amer. Math. Monthly}  56:672–676,  (1949).
%
%\bibitem{another_unique_key}
%Olga Taussky and John Todd.
%\newblock Cholesky, Toeplitz and the triangular factorization of
%symmetric matrices.
%\newblock {\em  Numer. Algorithms}, 41(2):197--202, 2006.
%\end{thebibliography}


\end{document}


% Please leave the next bit alone
%%% Local Variables:
%%% mode: latex
%%% TeX-master: "../ILAS-2022"
%%% End:
